\chapter{The Growing Tree}
\label{ch:growing_tree}

The spine of the guide is an eight-iteration search on our four-legal-move endgame position. Because there are only four legal moves at the root, every single node in the tree can be drawn explicitly at fixed coordinates without schematic shortcuts.

\section{The Eight Search Iterations}

We trace the search iteration by iteration. A FIG-2.x frame captures two moments: the \textbf{tree} displays the state \textit{after} iteration $x$ completes (new leaf evaluated, visit counts incremented, backpropagated values), while the \textbf{S-bar strip} below displays the scores compared \textit{before} iteration $x$'s pick --- the numbers that caused the selection. The tallest bar ($S = Q + U$) carries gold highlight and matches the candidate move marked with the gold selection ring in the tree above.

\begin{realdata}[Iteration 0]
\begin{center}
\input{figures/fig_2_0.tex}
\captionof{figure}{\textbf{FIG-2.0: Iteration 0 (Starting State).} All four children unvisited (dashed grey). All $Q$ values are equal to $Q_{\text{FPU}} = +0.97602$. Ranking is purely decided by policy priors.}
\label{fig:2_0}
\end{center}
\end{realdata}

\vspace{1ex}

\begin{realdata}[Iteration 1]
\begin{center}
\input{figures/fig_2_1.tex}
\captionof{figure}{\textbf{FIG-2.1: Iteration 1.} Kd6 selected ($S = 1.76358$). Expanded leaf returns $+0.96766$. Kd6 becomes solid green burst; backpropagation updates root to $Q = +0.97184$.}
\label{fig:2_1}
\end{center}
\end{realdata}
\clearpage

\begin{realdata}[Iteration 2]
\begin{center}
% GENERATED by tools/make_figures.py — do not edit.

\begin{tikzpicture}[scale=0.72, every node/.style={transform shape}]
  \node[vgroot] (root) at (0,0) {Root\\ \tiny $n{=}2$\quad $Q{=}+0.97655$};
  \vgheat{1}{1}\node[vgnode, fill=vgfill, inner sep=1.8pt] (kd6) at (-4.5,-1.8) {Kd6\\ \tiny $n{=}1$\ \ $Q{=}+0.96766$};
  \vgheat{1}{1}\node[vgnew, fill=vgfill, inner sep=1.8pt] (kf6) at (-1.5,-1.8) {Kf6\\ \tiny $n{=}1$\ \ $Q{=}+0.98598$};
  \node[vgunvisited, inner sep=1.8pt] (kf5) at (1.5,-1.8) {Kf5\\ \tiny $n{=}0$\ \ $Q_{\mathrm{FPU}}{=}0.66460$};
  \node[vgunvisited, inner sep=1.8pt] (kd5) at (4.5,-1.8) {Kd5\\ \tiny $n{=}0$\ \ $Q_{\mathrm{FPU}}{=}0.66460$};
  \vgedgewidth{1}{1}\draw[vgedge, line width=\vgEW pt] (root) -- (kd6);
  \vgedgewidth{1}{1}\draw[vgedge, line width=\vgEW pt] (root) -- (kf6);
  \vgedgewidth{0}{1}\draw[vgedge, line width=\vgEW pt] (root) -- (kf5);
  \vgedgewidth{0}{1}\draw[vgedge, line width=\vgEW pt] (root) -- (kd5);
  \draw[vgpath] (root) -- (kf6);
  \vgnewmark{kf6}
  \vgselectmark{kf6}
  \draw[vgbackprop] (kf6.north) to[out=40, in=320] node[midway, right=1mm, font=\tiny\bfseries, text=IdeaGreen] {$+0.98598$} (root.east);
  \vgdeltadown{kd6}{\tiny $U$ halved: $\div 2$}
  \vgdeltadown{kf5}{\tiny FPU drop}
  \vgdeltadown{kd5}{\tiny FPU drop}
  \begin{scope}[yshift=-9.5cm]
    \node[anchor=west, font=\tiny\itshape, text=WorkGrey] at (-6.0, 3.6) {Scores compared before this pick ($S = Q + U$):};
    \vgsbaraxis{-6.0}{6.0}{1.8}
    \vgsbarS{-4.5}{0.96766}{0.39380}{Kd6}{meas}
    \vgsbarS{-1.5}{0.75015}{0.77190}{Kf6}{sel}
    \vgsbarS{1.5}{0.75015}{0.09389}{Kf5}{fpu}
    \vgsbarS{4.5}{0.75015}{0.09180}{Kd5}{fpu}
  \end{scope}
\end{tikzpicture}

\captionof{figure}{\textbf{FIG-2.2: Iteration 2 (The PUCT Demonstration).} Kd6 has the best measured score ($+0.96766$ vs pre-pick FPU $0.75015$, grey bars) yet is \textbf{not} selected. Its $U$ halved ($1+0 \rightarrow 1+1$ denominator), allowing Kf6 ($S = 1.52205$) to take the lead. Post-pick FPU updates to $0.66460$ (dashed nodes).}
\label{fig:2_2}
\end{center}
\end{realdata}

\vspace{1ex}

\begin{realdata}[Iteration 3]
\begin{center}
\input{figures/fig_2_3.tex}
\captionof{figure}{\textbf{FIG-2.3: Iteration 3 (Selection Recurses).} Kf6 selected ($S = 1.53183$). Selection does not stop at root; it recurses to Kf6's child node Kf8 (depth 2). Leaf value $+0.95129$ backs up.}
\label{fig:2_3}
\end{center}
\end{realdata}
\clearpage

\begin{realdata}[Iteration 4]
\begin{center}
\input{figures/fig_2_4.tex}
\captionof{figure}{\textbf{FIG-2.4: Iteration 4.} Kd6 selected ($S = 1.64983$). Selection recurses down Kd6 branch to new leaf Kd8 (depth 2). Leaf returns $+0.99759$, updating Kd6's average $Q$ to $+0.98262$.}
\label{fig:2_4}
\end{center}
\end{realdata}

\vspace{1ex}

\begin{realdata}[Iteration 5]
\begin{center}
\input{figures/fig_2_5.tex}
\captionof{figure}{\textbf{FIG-2.5: Iteration 5.} Kd6 selected ($S = 1.50779$). Selection expands Kf7 under Kd6 (depth 2). Leaf returns $+0.99992$. Kd6 subtree widens.}
\label{fig:2_5}
\end{center}
\end{realdata}
\clearpage

\begin{realdata}[Iteration 6]
\begin{center}
% GENERATED by tools/make_figures.py — do not edit.

\begin{tikzpicture}[scale=0.72, every node/.style={transform shape}]
  \node[vgroot] (root) at (0,0) {Root\\ \tiny $n{=}6$\quad $Q{=}+0.97958$};
  \vgheat{3}{3}\node[vgnode, fill=vgfill, inner sep=1.8pt] (kd6) at (-4.5,-1.8) {Kd6\\ \tiny $n{=}3$\ \ $Q{=}+0.98839$};
  \vgheat{3}{3}\node[vgnode, fill=vgfill, inner sep=1.8pt] (kf6) at (-1.5,-1.8) {Kf6\\ \tiny $n{=}3$\ \ $Q{=}+0.97196$};
  \node[vgunvisited, inner sep=1.8pt] (kf5) at (1.5,-1.8) {Kf5\\ \tiny $n{=}0$\ \ $Q_{\mathrm{FPU}}{=}0.66763$};
  \node[vgunvisited, inner sep=1.8pt] (kd5) at (4.5,-1.8) {Kd5\\ \tiny $n{=}0$\ \ $Q_{\mathrm{FPU}}{=}0.66763$};
  \vgedgewidth{3}{3}\draw[vgedge, line width=\vgEW pt] (root) -- (kd6);
  \vgedgewidth{3}{3}\draw[vgedge, line width=\vgEW pt] (root) -- (kf6);
  \vgedgewidth{0}{3}\draw[vgedge, line width=\vgEW pt] (root) -- (kf5);
  \vgedgewidth{0}{3}\draw[vgedge, line width=\vgEW pt] (root) -- (kd5);
  \node[vgnode] (kf8) at (-1.0,-3.6) {Kf8\\ \tiny leaf $V{=}+0.95129$};
  \draw[vgedge] (kf6) -- (kf8);
  \node[vgnode] (kd8) at (-5.8,-3.6) {Kd8\\ \tiny leaf $V{=}+0.99759$};
  \draw[vgedge] (kd6) -- (kd8);
  \node[vgnode] (kf7) at (-3.4,-3.6) {Kf7\\ \tiny leaf $V{=}+0.99992$};
  \draw[vgedge] (kd6) -- (kf7);
  \node[vgnew] (e6_kf8) at (-1.0,-5.4) {e6\\ \tiny leaf $V{=}+0.97860$};
  \draw[vgedge] (kf8) -- (e6_kf8);
  \draw[vgpath] (root) -- (kf6);
  \draw[vgpath] (kf6) -- (kf8);
  \draw[vgpath] (kf8) -- (e6_kf8);
  \vgnewmark{e6_kf8}
  \vgselectmark{kf8}
  \vglane{e6_kf8}{root}{6.5}{-6.3}{$+0.97860$}
  \begin{scope}[yshift=-9.5cm]
    \node[anchor=west, font=\tiny\itshape, text=WorkGrey] at (-6.0, 3.6) {Scores compared before this pick ($S = Q + U$):};
    \vgsbaraxis{-6.0}{6.0}{1.8}
    \vgsbarS{-4.5}{0.98839}{0.44039}{Kd6}{meas}
    \vgsbarS{-1.5}{0.96863}{0.57547}{Kf6}{sel}
    \vgsbarS{1.5}{0.66779}{0.21000}{Kf5}{fpu}
    \vgsbarS{4.5}{0.66779}{0.20531}{Kd5}{fpu}
  \end{scope}
\end{tikzpicture}

\captionof{figure}{\textbf{FIG-2.6: Iteration 6.} Kf6 selected ($S = 1.54411$). Selection reaches depth 3 at e6 (under Kf8). Leaf returns $+0.97860$. Search deepens autonomously.}
\label{fig:2_6}
\end{center}
\end{realdata}

\vspace{1ex}

\begin{realdata}[Iteration 7]
\begin{center}
\input{figures/fig_2_7.tex}
\captionof{figure}{\textbf{FIG-2.7: Iteration 7.} Kd6 selected ($S = 1.47084$). Selection reaches depth 3 at e6 (under Kd8). Leaf returns $+0.97060$.}
\label{fig:2_7}
\end{center}
\end{realdata}
\clearpage

\begin{realdata}[Iteration 8]
\begin{center}
\input{figures/fig_2_8.tex}
\captionof{figure}{\textbf{FIG-2.8: Iteration 8 (Verification).} State after 7 backed-up values. Hand arithmetic matches \texttt{lc0.exe go nodes 8} visit counts and $Q$ values to 5 decimal places.}
\label{fig:2_8}
\end{center}
\end{realdata}

\vspace{1ex}

\begin{realdata}[Search Time-Lapse]
\begin{center}
\input{figures/fig_2_9.tex}
\captionof{figure}{\textbf{FIG-2.9: The whole search time-lapse.} Sequence of tree states across iterations 1 to 8 showing structural growth.}
\label{fig:2_9}
\end{center}
\end{realdata}
\clearpage

\section{The Sign Flip Convention}

Evaluations are recorded from the point of view of the player to move at that node. As values walk back up the tree, the sign flips at every ply:

\begin{realdata}[Sign Flips]
\begin{center}
\begin{minipage}{\linewidth}
\captionsetup{type=figure}
\centering
\subcaptionbox{Leaf Evaluation (White POV)}{\input{figures/fig_2_10a.tex}}\\[1ex]
\subcaptionbox{Black-to-move Edge}{\input{figures/fig_2_10b.tex}}\\[1ex]
\subcaptionbox{Root Edge (White POV)}{\input{figures/fig_2_10c.tex}}
\captionof{figure}{\textbf{FIG-2.10: Iteration 3 Sign Flips.} New leaf after 1.Kf6 Kf8 gives $+0.95129$ (White). Black edge receives $-0.95129$. Root edge receives $+0.95129$, updating Kf6 average to $0.968635$.}
\label{fig:2_10}
\end{minipage}
\end{center}
\end{realdata}

\section{Emergent Depth}

\begin{established}[Autonomous Search Expansion]
\begin{center}
\input{figures/fig_2_11.tex}
\captionof{figure}{\textbf{FIG-2.11: Depth Emergence.} Iterations 1--2 stop at depth 1; 3--5 reach depth 2; 6--7 reach depth 3. Depth is an emergent outcome of $U$ decay, not a pre-configured parameter.}
\label{fig:2_11}
\end{center}
\end{established}
